%This is my super simple Real Analysis Homework template

\documentclass{article}
\usepackage[utf8]{inputenc}
\usepackage[english]{babel}
\usepackage[]{amsthm} %lets us use \begin{proof}
\usepackage[]{amssymb} %gives us the character \varnothing

\title{Homework 2}
\author{Your Name}
\date\today
%This information doesn't actually show up on your document unless you use the maketitle command below

\begin{document}
\maketitle %This command prints the title based on information entered above

Read the code in the \texttt{main.tex} file to see how the document is formatted. ``Math mode'' in \LaTeX  can be delimited by dollar signs or by beginning and ending an \texttt{equation} or \texttt{aligned} environment. I suggest making a copy of the original \LaTeX file and editing the copy in case something breaks and you're not sure how to get back to the functioning file. As with any code, make incremental changes and compile them so that if you make a mistake it is easy to undo.

In order to generate a \texttt{.pdf} file, you will use the command line command

\begin{verbatim}
pdflatex file_name.tex
\end{verbatim}

You will need both \texttt{latex} and the \texttt{pdflatex} tools installed. Both should be included if you follow the instructions for your platform here:

\begin{verbatim}
https://www.latex-project.org/get/
\end{verbatim}




Let $m:\mathcal{A}\rightarrow [0,\infty)$ be a set function where $\mathcal{A}$ is a $\sigma$-algebra. Assume $m$ is countably additive over countable disjoint collections of sets in $\mathcal{A}$.
%Basically, you type whatever text you want and use the $ sign to enter "math mode".
%For fancy calligraphy letters, use \mathcal{}
%Special characters are their own commands

\subsection*{Problem 1}
Given sets $A$, $B$, and $C$, if $A\subset B \textrm{ and } B \subset C$, then $A \subset C$.
\begin{proof}
Other symbols you can use for set notation are
\begin{itemize}
\item$A \supset B \supseteq C \subset D \subseteq E$. Also $\varnothing \textrm{vs} \emptyset$
\item$\cup$ and $\cup_{k=1}^\infty E_k$
\item$\cap$ and $\cap_{x \in \mathbb{N}} \{\frac{1}{\sqrt[3]{x}}\}$
\item$\bigcup$ and $\bigcap\limits_{k=0}^n$ and $\bigcap$
\item most Greek letters $\sigma \pi \theta \lambda_i e^{i\pi}$
\item $\int_0^2 ln(2)x^2sin(x) dx$
\item$\leq < \geq > = \neq$
\end{itemize}
If you want centered math on its own line, you can use a slash and square bracket.\\
\[
\left \{
\sum\limits_{k=1}^\infty l(I_k):A\subseteq \bigcup_{k=1}^\infty \{I_k\}
\right \}
\]
The left and right commands make the brackets get as big as we need them to be.
\end{proof}

\clearpage %Gives us a page break before the next section. Optional.
\subsection*{Problem 2}
Given...
\begin{proof}
Let $\epsilon>0$.
If you have a shorter statement that you still want centered, use two \$\$ on either side.
$$\exists \textrm{ some } \delta>0 \mid ...$$
\end{proof}

\subsection*{Problem 3}
%
\begin{proof}
%
\end{proof}

\section*{Section 2.2}
%
\subsection*{Problem 6}
Blah
\subsection*{Problem 7}
Blah
\subsection*{Problem 10}
Blah

\end{document}
